\documentclass[11pt,a4paper]{article}
\usepackage[margin=22mm]{geometry}
\usepackage{amsmath,amssymb,amsthm,booktabs,microtype}
\usepackage[T1]{fontenc}
\usepackage{lmodern}
\usepackage[colorlinks=true,linkcolor=blue,urlcolor=blue]{hyperref}
\hypersetup{pdftitle={Subset sums: anchored hole bounds and the audited bound 91},pdfauthor={Sum-of-Subset research project}}
\newcommand{\FS}{\operatorname{FS}}
\newcommand{\Z}{\mathbb Z}
\newcommand{\T}{\mathcal T}
\newtheorem{theorem}{Theorem}[section]
\newtheorem{lemma}[theorem]{Lemma}
\newcommand{\RootGroups}{28,180}
\newcommand{\RootSystems}{271,156}
\newcommand{\CoarseGroups}{27,795}
\newcommand{\NestedGroups}{385}
\newcommand{\IntervalChecks}{28,412}
\newcommand{\ExtremumPoints}{74,021}
\newcommand{\TreeRoots}{1,342}
\newcommand{\BinaryNodes}{3,320}
\newcommand{\TreeLeaves}{4,662}
\newcommand{\MaxDepth}{14}
\newcommand{\FinalSystems}{274,476}
\newcommand{\Vectors}{4,130}
\newcommand{\MinGap}{3/2519}
\newcommand{\NegativeControls}{36}
\title{Subset Sums in a Fourfold Sumset\\\large Anchored hole bounds and integer center bases,\\and the audited bound $N\le91$}
\author{Sum-of-Subset research project}
\date{17 October 2026}
\begin{document}
\maketitle
\begin{abstract}
We continue the audited bound $N\le92$ by excluding cardinality 92.
The former uniform inverse margin becomes negative in fourteen
orientations. A progression/non-progression dichotomy repairs this
gate: non-progressions improve both the support and defect budgets;
two progression supports are treated exactly by residue classes.
The only unequal-step case left by the cardinality budget contradicts
the subsequent peeling budget. We then use the persistent-hole theorem
earlier, after core normalization but before propagating the lattice.
A full-prefix central interval lowers its density threshold by one;
a sharp non-progression anchor strengthens the defect recurrence.
For the size-92 model, these reduce the high-hole covering cap
from 199 to 150.
For covering, integer row operations on actual contributing centers
replace the former one-coordinate propagation. All three-empty modes
become integral without assuming contribution by the leading translates.
Four empty units now occur; all 10,880 modes are retained, with partial
integer information in the 30 exceptional modes. Finally, the second
core cutoff supplies stronger leading quantiles. A cardinality charge
after trimming one or both extremes strengthens the twofold-sumset
bound without classifying individual core sets. Low-hole mass transport
also retains the stronger support-excess floors from the hole witnesses.
Exact integer
certificates, a standard-library checker, and complete run records
accompany the report. The stronger archived bound 84 remains pending
complete independent replay.
\end{abstract}

\section{Result and inherited structure}
For a finite real set $A$, let $\FS(A)$ include the empty sum and let
$4A=A+A+A+A$ allow repeated summands. Write $N$ for the largest feasible
cardinality under $\FS(A)\subseteq4A$.
\begin{theorem}\label{main}
Every feasible set has at most 91 elements; hence $71\le N\le91$.
\end{theorem}
The unchanged lower construction is
$([-31,31]\cap\Z)\cup\{\pm93,\pm155,\pm217,\pm713\}$.
The audited foundation $N\le92$ is \cite{previous}. It suffices to
exclude size 92. Such a set contains zero, since adjoining zero to
one without it preserves feasibility and would give 93 elements.
No deletion principle for arbitrary subsets is assumed.

We inherit the corrected integralization and peeling lemmas
\cite{foundation}, the sharp four-extreme and joint-covering bounds
\cite{overlap}, the threshold-17 non-progression theorem \cite{sharp},
and the exact prefix inverse \cite{prefix}. For a positive $r$-set,
\[
 T_r=\frac{r(r+1)}2,\quad
 \mu(X)=|\FS(X)|-1-T_r,\quad
 \rho(X)=\mu(X)+2-2\delta_0(X),
\]
where $\delta_0$ counts $x$ for which $2x$ is not a distinct-element
two-sum of $X\cup\{0\}$. The maximum is defective, so $\rho\le\mu$.
The general floors are $\rho\ge-8$ and $\rho\ge0$ for $r\ge14$.
For non-progression support, $\mu\ge r-3$, $\rho\ge r-14$ for $r\ge6$,
and $\rho\ge r-5$ for $r\ge17$.

Peel as before to $D\supset D_1\supset E=P\cup\{0\}\cup(-Q)$, with
\[
 d=84,\quad e=76,\quad a+b=m=75,\quad I=(a-37)(a-38),
 \quad\Lambda=1445+I,\quad\Omega=500-I.
\]
Put $\rho=\rho(P)+\rho(Q)$ and $\mu=\mu(P)+\mu(Q)$.
The exact inherited identity is
\begin{equation}\label{identity}
 \rho+2\vartheta+c+\T=\Omega,\qquad
 \T=\kappa+\ell_1+\ell_2+g,
 \qquad c=2|2D|+18d+85-|\FS(D)|.
\end{equation}
All displayed slacks other than $\rho$ are nonnegative, and
$\kappa=|\FS(P)-\FS(Q)|-|\FS(P)|-|\FS(Q)|+1$.
The deficient-sign inequality is $1+T_{76}>24\cdot92+77$.
The paired bootstrap gives $q^2\le2593$ with odd $q$, hence $q\le49$
and $a,b\ge13$. At size 13 on one side, $\Omega=-100<-8$;
the other side has size 62 and nonnegative defect. This is impossible.
Both sides are now at least 14; $\Omega\ge0$ leaves $16\le a\le59$.

\section{Repairing the simultaneous inverse gate}
The sufficient inverse condition used in the foundation is
\begin{equation}\label{inverse}
 \kappa<\min(|\FS(P)|,|\FS(Q)|)-3.
\end{equation}
Let $r=\min(a,b)$, $t=\max(a,b)$, and
$J=T_r-2-\Omega$. Its minimum is now $-21$.
The fourteen orientations with $J\le0$ have $22\le r\le28$.

\begin{lemma}[Use non-progression before the inverse theorem]\label{nonapgate}
In each exceptional orientation, \eqref{inverse} holds whenever at
least one side-support is not a progression.
\end{lemma}
\begin{proof}
If the side of size $r$ is non-progression, its support grows by at
least $r-3$ above $1+T_r$ and its defect is at least $r-5$.
Since $t>r$, $T_t\ge T_r+r-3$. Thus the margin in \eqref{inverse}
is at least $J+2r-8$. If only the larger side is known to be
non-progression, its defect gives margin at least $J+t-5$.
Here we used $\kappa\le\Omega-\rho$, directly from \eqref{identity}.
The two exact minima over the fourteen orientations are 20 and 24.
\end{proof}
The exceptional sizes, with both orientations included, are as follows.
The last column is the residual budget in the double-step progression
case treated below.
\begin{center}\small
\begin{tabular}{rrrrrr}\toprule
$r$&$\Omega$&$J$&$J+2r-8$&$J+t-5$&$\Omega-T_r$\\\midrule
22&260&$-9$&27&39&7\\
23&290&$-16$&22&31&14\\
24&318&$-20$&20&26&18\\
25&344&$-21$&21&24&19\\
26&368&$-19$&25&25&17\\
27&390&$-14$&32&29&12\\
28&410&$-6$&42&36&4\\\bottomrule
\end{tabular}
\end{center}

\begin{lemma}[Two progression supports]\label{AP}
Let $U=p[0,u-1]_{\Z}$ and $V=q[0,v-1]_{\Z}$, with
$p,q\ge1$, $\gcd(p,q)=1$, and $u,v\ge2$. Then
\[
 |U+V|=\begin{cases}
 uv,&p\ge v\text{ or }q\ge u,\\
 pu+qv-pq,&p<v\text{ and }q<u.
 \end{cases}
\]
\end{lemma}
\begin{proof}
A collision requires index differences divisible by $q$ and $p$,
respectively. There are none in the first case. In the other case,
group indices of $U$ by $i=i_0+qk$, $0\le i_0<q$.
The corresponding sums in their residue class are translates of
$[0,v-1]$ by successive multiples of $p$. These intervals overlap
because $p<v$. If there are $K+1$ indices in the class, its union has
$v+pK$ values. Summing $K$ over the $q$ classes gives $u-q$,
and the asserted cardinality follows.
\end{proof}

Suppose now both supports are progressions in an exceptional orientation.
Integralize first, divide their steps by their gcd, and label the
side of size $r$ as $U$, so $u\ge1+T_r$ and $v\ge1+T_t$.
Reflection of the second support does not change the sumset cardinality.
The collision-free case gives $\kappa\ge T_rT_t>\Omega$.
In the other case,
\[
 \kappa=(p-1)(u-q)+(q-1)(v-1).
\]
Since $T_t>\Omega$, we must have $q=1$. Then
$\kappa=(p-1)(u-1)$, and $2T_r>\Omega$ forces $p\le2$.
If $p=1$, both supports already have a common step and no hull holes;
no inverse theorem is needed.

It remains to exclude $p=2,q=1$. The difference support is a full
interval because $v\ge2$. On this common lattice,
\[
 h_P+h_Q=\kappa=u-1,\qquad z=0,\qquad
 \Omega-\kappa\le\Omega-T_r\le19.
\]
This case has $h_U=|U|-1$, so we do \emph{not} apply any theorem
requiring a strict side-hole bound. Instead use the general span lemma
$\max X\le |X|+\rho(X)+h_X+2$, valid for every positive integer set
of size at least 15. For completeness, delete its largest element
$r+q$ and write $X'$ for the remaining set. If $\nu\ge0$ counts
destroyed diagonal defects, the exact increment identity is
$\rho(X)=\rho(X')+q-h_X+h_{X'}-2+2\nu$.
Since $|X'|\ge14$, all terms other than $q-h_X-2$ are nonnegative.
This proves the stated span bound without a primitivity hypothesis.
The two lattice-propagation comparisons below
apply directly, with $z=0$. They put $D_1,D$ on this lattice and
propagate zero holes. The sharp four-extreme theorem then gives
\[
 \T\ge\kappa+158,
\]
contradicting \eqref{identity} and $\Omega-\kappa\le19$.
The strict contradiction has margin at least 139.
This excludes every unequal-step progression case left by the budget.
Lemma~\ref{AP} is also checked against 8,428 directly enumerated
small progression pairs; its arbitrary-length proof is the residue
argument above.

\section{An earlier use of persistent holes}
In the remaining cases, the inverse theorem or equal progression
steps give a common core lattice with
\[
 h_P,h_Q\le\kappa,\quad z=h_P+h_Q-\kappa,
 \quad0\le z\le\min(h_P,h_Q),\quad h_X<|\FS(X)|-1.
\]
The inverse is not reapplied to the excluded double-step case.
The persistent-hole argument can be shortened and strengthened in two
places. Apply the central-interval theorem to the full prefix, avoiding
the deletion of its maximum. Then start the defect recurrence at the
sharp positive floor of the preceding non-progression prefix.

\begin{lemma}[Full-prefix endpoint holes]\label{fullprefix}
Let $Y\subseteq[1,i+\delta]_{\Z}$ have $i$ elements, where
$0\le\delta\le\beta$ and $i\ge2\beta+4$. With $S=\sum Y$,
\[
 \FS(Y)=[0,S]_{\Z}\setminus(H\cup(S-H)),\qquad
 H\subseteq[1,2\beta],\qquad |H|\le\beta.
\]
\end{lemma}
\begin{proof}
Theorem 1 of Lev \cite{lev} states that an $i$-element subset of
$[1,l]_{\Z}$ with $l\le3i/2-2$ has all subset sums in
$[2l-2i+1,S-(2l-2i+1)]_{\Z}$. Here $l=i+\delta$ meets that
hypothesis, so the interval $[2\delta+1,S-2\delta-1]$ is full.
Every missing lower sum in $[1,2\delta]$ is a missing member of $Y$,
because $2\delta\le i+\delta$; there are at most $\delta$ such
members. Complement symmetry determines the upper holes. Enlarging
the endpoint range from $2\delta$ to $2\beta$ proves the statement.
Unlike the previous counted-endpoint argument, this use does not need
to identify the number of deletions with the exact support increment.
\end{proof}

\begin{theorem}[Anchored persistent-hole bound]\label{holerho}
Let $X$ be a positive integer $r$-set with
$2\beta<h_X<|\FS(X)|-1$, where $\beta\ge2$. Choose
$\varepsilon\in\{0,1\}$ and
$r\ge L\ge\max(17+\varepsilon,2\beta+4)$. Then
\begin{equation}\label{persistent}
 \rho(X)\ge B_\varepsilon(r,\beta,L):=
 \min\left\{\begin{aligned}
 &\varepsilon(L-6)+(\beta-1)(r-L+1),\\
 &T_{L-\varepsilon}-r-1,\\
 &T_L+L-6\beta-r-7
 \end{aligned}\right\}.
\end{equation}
\end{theorem}
\begin{proof}
Assume $\rho_r$ is less than all three terms. The strict upper hole
condition forces $X$ to be primitive: a gcd at least two creates at
least $|\FS(X)|-1$ hull holes. Its support is non-progression because
a primitive progression support containing zero is a full interval.
For the increasing prefixes $X_i$, put
$y_i=|\FS(X_i)|-|\FS(X_{i-1})|-i$. The inherited increment identity
gives $\rho_i\ge\rho_{i-1}+y_i-2$; for $i\ge6$, $y_i=0$
forces progression support.

If the last progression prefix has size $t\ge L-\varepsilon$, its
next support translate is disjoint. Indeed, an on-lattice translate
that meets or abuts the previous progression preserves progression
support; an off-lattice translate is disjoint. Thus
$y_{t+1}\ge T_t-t$, and all subsequent increments are at least one.
Using $\rho_t\ge0$ gives
\[
 \rho_r\ge T_t-t-2-(r-t-1)=T_t-r-1
 \ge T_{L-\varepsilon}-r-1,
\]
a contradiction. All prefixes from $L-\varepsilon$ onward are
therefore non-progressions. If $\varepsilon=1$, the sharp theorem at
size $L-1\ge17$ gives $\rho_{L-1}\ge L-6$; if
$\varepsilon=0$, the general floor gives $\rho_{L-1}\ge0$.
If $y_i>\beta$ throughout $L\le i\le r$, telescoping from this
anchor gives the first term of \eqref{persistent}, again impossible.
Choose $i\ge L$ with $1\le y_i\le\beta$.

Since $i\ge2\beta+4\ge\beta+5$, the inherited small-increment
inverse applies and writes $X_i=tY$ with $Y$
primitive, $|Y|=i$, and $\max Y\le i+\beta$.
The scale $t$ is an integer since $X_i$ is integral and $\gcd(Y)=1$.
If $t>1$, final primitivity forces a first later element off $t\Z$,
say at index $k$. Its support translate is disjoint. The same increment
calculation, with $\rho_{k-1}\ge0$, gives
$\rho_r\ge T_{k-1}-r-1\ge T_L-r-1$, contradicting the second term.
Hence $t=1$ and Lemma~\ref{fullprefix} applies to $X_i$.

Suppose the endpoint-hole description holds at size $v\ge i$, with
sum $S$, and let $x$ be the next element. If $x\ge S-4\beta$,
its translate adds at least $S-6\beta$ points above $S$.
For $x\le S$, at least $x-2\beta$ of the last $x$ positions are
present; for $x>S$, the whole translate is new. Consequently
$y_{v+1}\ge S-6\beta-v-1$. Using the sharp floor
$\rho_v\ge v-5$ and all later increments at least one gives
\[
 \rho_r\ge S+v-6\beta-r-7\ge T_L+L-6\beta-r-7,
\]
contradicting the third term. Therefore $x<S-4\beta$.
Also $x>2\beta$ since $x$ exceeds the maximum of a prefix of size
at least $2\beta+4$. The two central intervals overlap or abut;
the lower holes remain $H$, and complement symmetry supplies the upper
holes. This propagates to $X_r$, giving $h_X\le2\beta$, the final
contradiction.
\end{proof}
For $z\ge19$, $(\beta,L,\varepsilon)=(9,22,1)$ gives respective
floors $136,152,168$ at sizes $36,38,40$, instead of the previous
$112,128,144$. The new helper \path{verification/nonap_hole_rho_full.py}
checks 22,335 parameter rows, retaining both values of $\varepsilon$.
It also checks the endpoint description directly in 88,392 deleted
intervals at the boundary $i=2\beta+4$, for $2\le\beta\le6$ and
$0\le\delta\le\beta$. These finite checks supplement the general
proof; they do not substitute for the central-interval theorem.

Let $\mathcal R(r,z)$ be the maximum of zero and these accepted floors
with $2\beta<z$, $\beta\le20$, both permitted values of
$\varepsilon$, and all permitted $L$. Let $F(r,u)$ be the unchanged exact
prefix inverse of \cite{prefix}, and write
$W(r,u,z)=\max\{F(r,u),\mathcal R(r,z)\}$.

Put $C=2m+4e+14-\Lambda=-977-I$. The inherited inner comparisons are
\[
 |\FS(D_1)|-2|\FS(E)|\le
 \min\{C+\Omega+\rho-\mu+2z,\ C+3\Omega-\rho-\mu\}.
\]
For $z\le14$, $\rho\le\mu$ gives the upper bound $-449-2I$.
For $z\ge15$, first use $\mu\ge69$ and the applicable non-progression
floors. The only unresolved sizes have $a,b\ge28$.
The inherited effective-hole witnesses at $\beta=7$ supply lower
bounds $u_r=r+C_7(r)+1$ for $\mu(X)$; now use
\[
 \rho+\mu\ge u_a+u_b+W(a,u_a,15)+W(b,u_b,15).
\]
The high-hole inner comparison is at most $-31$ over all 44 orientations.
This use of \eqref{persistent} is legitimate because $P,Q$ have already
been normalized, even though the peeled layers are not yet known to
lie on their lattice. The outer comparison is
\[
 \Lambda-(4d-10+\Omega)=619+2I>0.
\]
Disjoint translates would violate these strict factor-two comparisons,
so $D_1,D$ lie on the core lattice. In the double-step progression case
of the previous section only the low-hole comparison is used; its
derivation needs the general span lemma, not a primitive side support.

Hole propagation and the sharp four-extreme theorem now give
\begin{equation}\label{peeling}
 w\le4z,\quad \T\ge\kappa+158-3z,
 \qquad\T\ge\max(53,z,158-2z).
\end{equation}
Neither size-16 boundary can survive these budgets. Thus all subsequent
calls to $F$ and $\mathcal R$ are within their proved range $r\ge17$.

\section{Integer bases of contributing centers}
Write the original four positive outer magnitudes as
$p_1>p_2>p_3>p_4>M_p$, and similarly use $q_i>M_q$ on the negative
side. The core lattice already contains $M_p,M_q$ and all of $D$.
Set
\[
 c_p=p_1-p_2-p_3-p_4,\qquad
 b_p=-p_1+2p_2-p_3-p_4,
\]
and reflect for $q$. The twenty ordinary center forms are
\begin{align*}
 v_{pi}&=c_p+p_{i+1},&v_{p,i+4}&=c_p-q_{i+1}&&(0\le i<4),\\
 v_{p8}&=p_2-p_3-p_4,&v_{p9}&=b_p,
\end{align*}
and their reflected forms. The eighteen-unit selection and the two
branches $b_p+b_q\le0$ and $b_p+b_q>0$ are inherited from
\cite{overlap,previous}. Every actual contributing ordinary center is
integral: its translate of the integral set $D$ intersects the integral
target. The actual member of a split pair is treated in the same way.

\begin{lemma}[Integer row certificates]\label{integerbasis}
Let $M$ have, as rows, the coefficient vectors of known integral center
forms in the eight outer coordinates. If an integer row vector $v$ is
an integer combination of rows of $M$, then $v\cdot(p,q)$ is integral.
In particular, an integer matrix $U$ with $UM=I_8$ proves integrality
of all eight outer coordinates.
\end{lemma}
\begin{proof}
Integer linear combinations of integers are integers. This also proves
the claim for individual gaps and for the branch form $b_p+b_q$.
\end{proof}
The new helper \path{verification/integral_centers.py} uses row swaps,
sign changes, and integer row subtraction to construct an echelon
basis by the Euclidean algorithm. It tracks the integer combinations
and checks $H=UM$ independently by matrix multiplication. A target
form is accepted only when integer elimination against $H$ leaves zero
remainder. A nonunit pivot is never divided away. Thus row reduction
is a means of finding explicit integer combinations, not a rational
rank test.

For zero through three empty units, the ordinary contributing forms
alone recover all coordinates in respectively $4,70,578,2992$ modes.
Leading contribution is not needed for these assertions. For example,
the old exception in branch B with missing $p0,p1,p8$ and second-pair
choice $p$ has
\[
 p_2=c_p+v_{p9}+2p_3+2p_4.
\]
The retained $q$ centers recover its side, $v_{p4}+q_1$ recovers $c_p$,
and the retained own centers recover $p_3,p_4$. The displayed integer
combination then recovers $p_2$, followed by $p_1$.
This is the combination missed by the former one-unknown rule.

Four empty units are now possible. All 10,880 modes are enumerated.
Ordinary centers recover 10,812 of them. If the residual cap is below
$2d-1=167$, a leading center may be added \emph{only on a side whose
own-center block is nonempty}. That block charges the actual leading
loss, which is then less than $|2D|$ and proves contribution.
This recovers 38 more modes, for a total of 10,850.
The remaining 30 modes are not discarded: 12 have rank seven, and 18
have rank eight with a nonunit obstruction. The checker retains every
individually proved integral gap or branch form and leaves other strict
gaps weak in the real relaxation. Their integer-gap counts are
\[
 6\text{ modes with }4,\quad18\text{ with }5,\quad
 4\text{ with }6,\quad2\text{ with }7.
\]
No leading quantile is imposed on a side with no charged own-center block.

\section{A deeper leading quantile}
Retain $C_p$, the second-largest element of the positive core $P$, and
$C_q$ on the other side. The mass and independent-cutoff rows of
\cite{previous} remain valid. With
$x_{-1}=C_p,x_0=R_p,x_4=M_p$, and similarly $y_i$, these include
\begin{gather*}
 C_p\ge a-1,\quad C_p+1\le R_p,\quad
 C_p+R_p+x_1+x_2+x_3+M_p-S_p\le-T_{a-2},\\
 K:=|2D|\le M_p+M_q+x_i+y_h+1+T_{4-i}+T_{4-h}
 \qquad(-1\le i,h\le4).
\end{gather*}
The 36 tail bounds are inherited, not claimed as a new theorem here.
The new use of $C_p,C_q$ concerns the leading translates.

\begin{lemma}[Leading quantiles from a smaller core]\label{deep}
Let a leading center $c$ target $[1,S]_{\Z}$ and let $L$ be its actual
total loss from $c+2D$. If $L<2d-11$, then
\[
 c+2C_p+L\ge2d-11,\qquad
 c-S-2C_q-L\le12-2d.
\]
For $d=84$, the two right sides are 157 and $-156$.
\end{lemma}
\begin{proof}
Remove the five largest positive elements of $D$. The remaining set
has $d-5$ elements, maximum $C_p$, and a twofold sumset of at least
$2d-11$ points. Its translate is a subset of $c+2D$. Since its total
loss is at most $L<2d-11$, it contributes to $[1,S]$.
At most $c+2C_p$ positive integers lie at or below its maximum;
at most $L$ further points lie below 1. This proves the first inequality.
Remove the five most negative elements instead. The minimum is
$-C_q$. Counting integers from $c-2C_q$ to $S$, with at most $L$
points above $S$, proves the second inequality.
\end{proof}
The proof uses the same total loss in two alternative necessary
inequalities; it does not add the losses twice. In the model the rows
are enabled only when the entire group's residual cap is less than
157 and that side's leading loss is charged. They are used only in the
nested model, where $C_p,C_q$ are retained.

\section{Trim first, then bound the sumset}
The independent hull cutoffs can overcount the contribution of a
widely separated extreme. Its number of possible partners gives a
second, independent upper bound.
\begin{lemma}[Cardinality cost of trimmed extremes]\label{trim}
Let $D$ have $d$ elements and $F=D\setminus T$, where $|T|=s$.
Then
\[
 |2D|\le |2F|+sd-T_{s-1}.
\]
\end{lemma}
\begin{proof}
Use $2D=2F\cup(F+T)\cup2T$, with
$|F+T|\le s(d-s)$ and $|2T|\le s(s+1)/2$.
The sets may overlap; summing their cardinalities is an upper bound.
The last two terms sum to $sd-T_{s-1}$.
\end{proof}
Remove the $k$ largest positive and the $\ell$ most negative elements,
where $(k,\ell)\in\{(1,0),(0,1),(1,1)\}$. Put
$A=4-k$, $B=4-\ell$, and $s=k+\ell$.
The remaining extremes are $x_A,-y_B$ in the preceding notation.
The same independent-cutoff proof applied to $F$ gives, for
$-1\le i\le A$ and $-1\le h\le B$,
\begin{equation}\label{trimrow}
 K\le x_A+y_B+x_i+y_h+1+T_{A-i}+T_{B-h}+84s-T_{s-1}.
\end{equation}
There are $30+30+25=85$ such rows. They do not require the original
outer coordinates to be integral, and are valid for every actual $D$.
All mass-transport coefficients in these rows are zero. They are
retained in nested models, while coarse models may discard them.

For illustration, a size-92 relaxation in the four-empty family had
$R_p=41$, $R_q=38$, $x_3=50$, $y_3=47$ and allowed $K=373$.
Removing both extremes and using $i=h=0$ gives instead
\[
 K\le50+47+41+38+1+6+6+167=356.
\]
This is an excluded relaxation point, not a feasible construction.
The proof no longer needs an enumeration of small-defect core sets
to remove this obstruction.

The checker directly enumerates the same 784 small set pairs used for
the inherited core-cutoff audit. All 66,640 trimmed cutoff checks pass;
735 of these pairs have a strictly better minimum bound than the 36
untrimmed cutoffs. This finite check supplements Lemma~\ref{trim} and
the general cutoff proof. The three-element example $D=\{-3,0,1\}$
shows why the two-extreme increment cannot be reduced from $2d-1$
to $2d-2$: $|2D|=6$ and $|2\{0\}|=1$.

\section{Complete size-92 model universe}
Let $U(r,z)$ be the inherited effective-hole lower bound for $\mu$,
with $U(r,1)=U(r,2)=r-3$. For positive holes the scalar cap is
\[
 c\le500-I-W(a,U(a,z),z)-W(b,U(b,z),z)-\max(53,z,158-2z).
\]
For zero holes use $c\le342-I$. In particular $342\ge4d=336$ but
$342<5d$; retaining only three empty units would be incomplete.
There are 460 surviving low-hole states and 2,968 high-hole states.
The positive low bands are
\begin{center}
\begin{tabular}{lrrrr}\toprule
Band&Hole range&Combined defect floor&$c+I$ cap&Endpoint\\\midrule
$P$&1--2&65&281&2\\
$Q$&3--4&77&273&4\\
$R$&5--14&100&270&14\\\bottomrule
\end{tabular}
\end{center}
The latter two defect floors follow from the exact $\beta=1,2$ witnesses
and the prefix inverse, over all 42 orientations with $17\le a\le58$.
We also retain the support excess supplied by those same witnesses.
For every $17\le r\le58$, the effective-hole parameter checks accept
$(\beta,C)=(1,5)$ and $(2,r-6)$, respectively. Since the conclusion is
$\mu\ge r+C+1$ when $h_X>2\beta$, the necessary mass floors are
\[
 \epsilon_r=\begin{cases}
 r-3,&1\le z\le2,\\
 r+6,&3\le z\le4,\\
 2r-5,&5\le z\le14.
 \end{cases}
\]
All 84 parameter witnesses are checked exactly. The latter two rows
use $h_X\ge z$ and the strict side-hole bound established after the
inverse gate. In particular, they are not applied to the excluded
double-step progression case. The mass inequalities become
\[
 S_p\ge T_a+\epsilon_a+z+x_1+x_2+x_3+M_p,
 \qquad 2S_p\ge2T_a+2\epsilon_a+2z+2M_p+6R_p+12,
\]
and their reflections. Keeping only the weaker $\epsilon_r=r-3$
in all three bands admits a residual $R_2$ relaxation; the stronger
floors preserve information already proved by the hole witnesses.
Each band now permits three empty units. Before endpoint transport
impose its constant cap; then move $z$ to its upper endpoint $Z$ and
increase $S_p,S_q$ by $Z-z$. The new mass floors depend only on side
size, so this transport remains valid. Every retained row has
$v_{S_p}+v_{S_q}+v_z\le0$, except the upper-hole row, which becomes
an equality. This invariant is checked for every rebuilt model.

For $z\ge15$, the strengthened bound gives $c\le150<2d$.
Thus two empty units are impossible, every charged leading loss is
below 157, and a one-empty model has residual cap at most
$150-84=66<80$. All high groups consequently license the deeper
leading quantiles; their one-empty cases also license core quantiles.
The actual contributor fixes the direction of the active-center test,
as in the preceding reports. The remaining states comprise 1,068
zero-empty and 1,022 one-empty states, counting each permitted empty
count separately. They are grouped by
$(a,j,U(a,z),U(b,z),v_p,v_q)$, retaining one common hole coordinate
across each group's interval. There are 54 groups, with
871,618 source-to-group row checks over 2,090 retained states.
With all other conditions fixed, replacing only the previous
persistent-hole floor reduces the complete structural universe from
35,372 to 28,180 groups; the high-hole portion drops from 9,916 to
2,724. The comparison is recorded in \path{hole_bound_comparison.json}
in the results directory. The former universe is only a comparison
of necessary conditions, not an additional completed certificate.
For $E_*=500-I-v_p-v_q$ and $\delta=c-84j$, retain together
\[
 \delta\le C_*-84j,\quad\delta+z\le E_*-84j,\quad
 \delta\le E_*-53-84j,\quad\delta-2z\le E_*-158-84j.
\]
The hole variable may not be maximized independently of the cap.

The mass rows are the preceding report's rows with $a+b=75$; the span
and cardinality budgets become
\[
 R_p+R_q+\delta-4z\le421-I-84j,\qquad
 S_p+S_q+\delta-2K-4z\le1596-84j.
\]
The dense-span lemma $2R\ge2r+\min(z,r-3)$ is unchanged, including
its valid chord on a high group's fixed hole interval. Joint block
deficits use the core size $e=76$ and are charged only once.
For a block of $k$ retained ordinary centers,
\[
 O\ge76k-(R_p+R_q+v_{\max}-v_{\min}+1)+\max_v(l_v+u_v).
\]
The inherited leading-plus-own-center block inequalities use the same
replacement of 77 by 76; no leading loss is charged for an absent block.
The coarse bound is $K\le M_p+M_q+R_p+R_q+21$; nested models retain
the five envelope levels, all 36 core cutoffs, the 85 trimmed cutoffs,
and six inner-loss positions.
The group-wide residual cap $C_g$ controls the quantile rows:
\begin{center}
\begin{tabular}{ll}\toprule
Rows&Required cap\\\midrule
Ordinary quantiles&$C_g<r_{\min}+5$\\
Core quantiles&$C_g<80$\\
Basic leading quantiles&$C_g<167$ and a charged own block\\
Nested leading quantiles&$C_g<159$ and a charged own block\\
Second-core leading quantiles&$C_g<157$ and a charged own block\\\bottomrule
\end{tabular}
\end{center}
Global quantiles use 84 and 83, core quantiles 80 and 79, and ordinary
quantiles $b+5$ and $a+4$ on side $p$. A leading nested level $i$
uses $159+2i$ and $158+2i$ for $0\le i\le4$.
Strict outer gaps and the strict branch-B row are licensed by their
individual integer row certificates, not by an empty-count shortcut.
All inherited binary alternatives remain available, with both children
required and no restriction on how many large-loss centers occur.

\section{Exact certificates and reproduction}
Every root, matrix and right side is rebuilt from these necessary
conditions. The low universe contains 268,432 integer root systems.
The complete verified statistics are
\begin{center}\small
\begin{tabular}{lr}\toprule
Quantity&Count\\\midrule
Structural groups&\RootGroups\\
Integer root systems&\RootSystems\\
Coarse / nested groups&\CoarseGroups\ / \NestedGroups\\
Symbolic interval checks&\IntervalChecks\\
Integer extremum evaluations&\ExtremumPoints\\
Integer roots requiring trees&\TreeRoots\\
Binary internal nodes / leaves&\BinaryNodes\ / \TreeLeaves\\
Maximum tree depth&\MaxDepth\\
Final integer leaf systems&\FinalSystems\\
Distinct integer vectors&\Vectors\\\bottomrule
\end{tabular}
\end{center}
For $Ax\le b$ and $x\ge0$, an integer multiplier $w\ge0$ with
$w^TA\ge0$ and $w^Tb<0$ is a contradiction. The minimum normalized
strict margin is $\MinGap$. Symbolic right sides are exact quadratics;
their integer extrema are checked at endpoints and around the vertex.
The complete integer parameter range is covered exactly once, and
every binary tree certifies both children. High-hole certificates
cover their common real hole interval simultaneously.
The saved multiplier inventory confirms use of the trimmed-sumset
rows in 2,787 vectors covering 324 structural groups, and use of the
second-core leading quantiles in 39 vectors covering 28 groups.
These counts are recorded in \path{method_usage.json}.

The \NegativeControls\ negative controls include omission of the
new $Z_4,Q_3,R_3$ families, an unjustified raw inverse gate, division
by a nonunit pivot, universal four-empty integrality, a leading
quantile without a charged block, and an attempt to exclude size 91.
The general mathematical reductions are proved here and in the cited
foundation; these finite checks are not a proof-assistant formalization.

Only Python 3.10+ and its standard library are required for verification:
\begin{verbatim}
python -I -S -B verification/upper91.py --output .work/upper91.json
python -I -S -B verification/upper91.py --exhaustive
\end{verbatim}
The data are \path{data/upper91_certificates.json}. The optional
\path{experiments/generate_upper91_certificates.py} uses SciPy to
discover candidates; exact integer arithmetic determines acceptance.
Standard and optimized Python outputs agree byte for byte. The exhaustive
result agrees except for its mode label, regenerated certificates are
byte-identical, and the unchanged upper-92 baseline is reproduced.
The large standard result is archived as \path{result.json.gz};
lossless decompression recovers the exact verifier-output bytes.
\path{result_archive.json} records both hashes and the roundtrip check.
Run receipts, hashes, center-basis witnesses, scalar data and logs are in
\path{results/2026-10-17/upper91/}. The Windows entry point runs all eleven
audited finite continuations. The archived endpoint 84 has not thereby
received a complete independent replay.

\begin{thebibliography}{9}
\bibitem{lev} V. F. Lev, \emph{On consecutive subset sums},
Discrete Mathematics 187 (1998), 151--160, Theorem 1.
\href{https://doi.org/10.1016/S0012-365X(98)80006-X}{doi:10.1016/S0012-365X(98)80006-X}.
\bibitem{previous} Sum-of-Subset research project,
\emph{Core cutoffs and the audited bound $N\le92$},
\path{reports/upper92_from93_20261016.tex}, 16 October 2026.
\bibitem{prefix} Sum-of-Subset research project,
\emph{An integer prefix inverse and the audited bound $N\le93$},
\path{reports/upper93_from94_20261016.tex}, 16 October 2026.
\bibitem{sharp} Sum-of-Subset research project,
\emph{A sharper non-progression lemma and the audited bound $N\le94$},
\path{reports/upper94_from95_20261016.tex}, 16 October 2026.
\bibitem{overlap} Sum-of-Subset research project,
\emph{Joint covering losses and the audited bound $N\le95$},
\path{reports/upper95_from96_20261016.tex}, 16 October 2026.
\bibitem{foundation} Sum-of-Subset research project,
\emph{The audited continuation from 115 to 103},
\path{reports/upper103_from115_20261014.tex}, 14 October 2026.
\end{thebibliography}
\end{document}
